% !PRINTED 13
\weberpage{13}

If one then sets
\weq{9}{z^2 = y,}
there results
\weq{10}{\frac{2\,dz}{\sqrt{(z^2-\alpha)(z^2-\beta)}} = \frac{dy}{\sqrt{y(y-\alpha)(y-\beta)}}}
and one has thus obtained, by means of the \gesperrt{quadratic substitution} (9), an elliptic differential under whose root sign there stands a function of the third degree with real roots. By \S\ 3 this differential can be brought by a linear substitution into the normal form
\[
\frac{d\zeta}{\sqrt{\zeta(1-\zeta)(1-\varkappa^2\zeta)}}
\]
and therein $\zeta$ and $\varkappa$ may be assumed to be positive proper fractions.

The \gesperrt{Legendre} normal form results from this if one sets:
\[
\zeta = \sin^2\varphi, \qquad d\zeta = 2\sin\varphi\cos\varphi\,d\varphi
\]
\weq{11}{\frac{d\zeta}{\sqrt{\zeta(1-\zeta)(1-\varkappa^2\zeta)}} = \frac{2\,d\varphi}{\sqrt{1-\varkappa^2\sin^2\varphi}}.}

\begin{center}
\S\ 5. The \gesperrt{Weierstrass} Normal Form.
\end{center}

\gesperrt{Weierstrass} made another normal form of the elliptic differential the basis of his investigations, namely the form
\weq{1}{du = \frac{dz}{\sqrt{4z^3 - g_2 z - g_3}},}
in which $g_2$, $g_3$ are constants that are called the first and second invariant of the differential. The general elliptic differential
\weq{2}{\frac{dx}{\sqrt{f(x)}},}
in which
\weq{3}{f(x) = a_0 x^4 + a_1 x^3 + a_2 x^2 + a_3 x + a_4}
holds, can be brought into the Weierstrass normal form by a \gesperrt{linear transformation}, provided the roots $x_1$, $x_2$, $x_3$, $x_4$ of $f(x)$ are known. This is done most simply in the following way.
